\section{Two-sided continuation geometry}\label{sec:dual}
The preceding lower does not require the suffix to be independent of the prefix. Its strength nevertheless depends on finding a common sublaw. We now separate three questions: the exact one-cut compression problem, the production of nonzero sublaws from raw kernels, and recursive realization of their geometry. This leads to a quantitative converse and construction on a class with correlated reports. No asymptotic coding limit is taken.

\subsection{The exact cut problem and common channels}
Keep the fixed physical exploration and bounded task of Section~\ref{sec:problem}. Write $\mu=\law(H)$ and let $K(du\mid h)$ be the actual conditional law of the complete suffix. It includes every public action, phase, cost and report available before the decision. Public exploration randomness is part of the visible history. Private physical/controller randomness is integrated into $K$ unless actually reported; it is never made free decoder information.

Let $\mathcal D_M(\Lambda,z)$ be the least excess over $B_T$ when computation before and after one cut is unrestricted, but the only prefix-dependent object crossing it is a label in $[M]$. The decoder then receives the whole physical suffix. This is a relaxation of online prediction, not an assertion that a one-cut code can be run at every cut. Let $D\subset\R^q$ be a fixed compact convex set containing all task means. Decoder functions below are Borel and $D$-valued.

\begin{theorem}[One-cut comparison and conditional continuation]\label{thm:cutduality}
For any prepared standard-Borel experiment with encoder-independent exploration,
\begin{equation}\label{eq:exactfunctional}
 \mathcal D_M(\Lambda,z)=\inf_{a_1,\ldots,a_M}
 \int\min_{i\le M}\int\norm{z(h,u)-a_i(u)}^2K(du\mid h)\,\mu(dh).
\end{equation}
If finite measures $\mu_-,\mu_+$ and a suffix probability $\eta$ satisfy
\[
 \mu_-\otimes\eta\le\Lambda\le\mu_+\otimes\eta,
\]
then, for $Z(h)=z(h,\cdot)$ in $L^2(\eta;\R^q)$,
\begin{equation}\label{eq:cutsandwich}
 q_M(Z_*\mu_-)\le\mathcal D_M(\Lambda,z)
 \le q_M(Z_*\mu_+).
\end{equation}
In particular, $c\mu\otimes\eta\le\Lambda\le C\mu\otimes\eta$ gives comparison with $q_M(Z_*\mu)$ within factors $c,C$, for every $M$.

There is also a conditional lower when global product domination vanishes. Let $W$ be a standard-Borel physical variable, measurable from $H$, which is genuinely supplied again in $U$, or is explicitly given to the decoder as a relaxation. Write $\rho=\law(W)$ and disintegrate the actual joint law as $\Lambda_w$. Suppose jointly measurable kernels $\mu_{-,w}$ and $\eta_w$ satisfy $\mu_{-,w}\otimes\eta_w\le\Lambda_w$ for $\rho$-almost every $w$. Then
\begin{equation}\label{eq:conditionalcut}
 R_\on(T,M)-B_T\ge
 \int q_M\bigl((h\mapsto z(h,\cdot))_*\mu_{-,w};L^2(\eta_w)\bigr)\,\rho(dw).
\end{equation}
No normalization of the finite measures is intended. In particular $q_M(t\nu)=tq_M(\nu)$ for $t\ge0$. Conditioning in \eqref{eq:conditionalcut} is not permission to leave a continuous register uncharged in an upper construction.
\end{theorem}
\begin{proof}
Projection subtracts $B_T$. Fix the independent algorithm seed, including all future private decoder coins. Once its cut label is fixed, the entire subsequent finite sequence of Borel updates is a single Borel function of the complete suffix. Thus there are at most $M$ decoder functions. For a fixed list, integration against $K$ gives $M$ Borel distortions as functions of $h$; selecting the first minimum is Borel and optimal. Infimizing over lists proves \eqref{eq:exactfunctional}. Averaging independent seeds cannot improve its deterministic infimum. Projecting into $D$ cannot increase loss.

The lower comparison follows by integrating against the dominated product and then minimizing in $L^2(\eta)$. For the upper choose an approximating Hilbert center list for $Z_*\mu_+$, project its values into $D$, choose Borel representatives, and assign each prefix to its nearest center in that Hilbert norm. The actual distortion is bounded above by the dominating product distortion. Let the approximation error tend to zero. Neither comparison optimizes a policy-dependent law.

For the conditional statement, a code remains an $M$-function code after conditioning on $W=w$ and independent seeds. Apply the lower argument in each fiber and integrate. To verify measurability without assuming a measurable family of optimizers, fix a countable algebra generating the suffix Borel sigma field and a countable dense subset of $D$. The bounded simple functions on this algebra with those values are countable and dense in $L^2(\eta_w;D)$ for every $w$. Hence each fiber quantization value equals an infimum over countably many $M$-tuples of such functions. Each associated integral is a Borel function of $w$ by the kernel integration theorem. Their countable infimum is measurable. Conditional versions are immaterial because every dominated product is absolutely continuous with respect to its actual fiber law.
\end{proof}

Formula~\eqref{eq:exactfunctional} is the one-shot functional source-coding problem with decoder side information, written for the actual task conditional mean. We do not claim that ordinary encoder/decoder reduction as new. The useful implications are its comparison with acquired Hilbert measures, its conditional localization, and the recursive realization theorem below. Side information is genuinely causal here because it arrives after the retained-state cut; its coding problem after fixing the cut is an ordinary one-way problem.

\begin{lemma}[Raw constructions of common continuation laws]\label{lem:minorization}
The following give certificates for Theorem~\ref{thm:cutduality}.

\textup{(i)} If $K(du\mid h)$ has density $k(h,u)$ relative to a probability $\eta$ and $k(h,u)\ge\alpha(h)\ge0$ for $\mu\otimes\eta$-almost every $(h,u)$, then $\mu_-=\alpha\mu$. For a finite suffix alphabet with $\eta(u)>0$, one may take $\alpha(h)=\min_u K(u\mid h)/\eta(u)$. Upper density bounds give the upper comparison in the same way.

\textup{(ii)} Suppose a raw post-cut block, including \emph{all} its visible reports, costs and phases $R$, and its post-block hidden state $X'$, has conditional kernel
\begin{equation}\label{eq:regeneration}
 L(d r,dx'\mid h)\ge\alpha(h)\kappa(d r)\pi(dx').
\end{equation}
Suppose further that the subsequent controlled continuation has a kernel $F(du'\mid r,x')$ independent of the earlier prefix. This requires resetting or retaining in $x'$ every controller variable used later. Then the complete suffix $U=(R,U')$ has common sublaw
\[
 K(du\mid h)\ge\alpha(h)\eta(du),\qquad
 \eta(dr,du')=\kappa(dr)\int F(du'\mid r,x')\pi(dx').
\]
The same constructions hold after disintegration over $W$.
\end{lemma}
\begin{proof}
Multiply the density inequality by $\mu(dh)$ for (i). For (ii), integrate \eqref{eq:regeneration} against the nonnegative continuation kernel $F$; its normalization makes $\eta$ a probability. A hidden-state reset alone is insufficient if a prefix-dependent controller or an earlier visible block report is omitted. The displayed joint block minorization is precisely what prevents that omission. Integrating each fiber proves the final assertion.
\end{proof}

In (ii) a Doeblin block of a prepared instrument produces the product certificate before any geometric estimate is made. Its mass is $\int\alpha\,d\mu$, and this mass multiplies the squared-distortion lower. Conditioning on a rare successful block without its mass would be a different experiment.

\begin{example}[A diagonal interface with no global product certificate]\label{ex:diagonal}
Prepare an independent uniform $W\in[0,1]$ and $d$ fair bits $X_1,\ldots,X_d$. Report $W$ and then the bits by paid calls. After the cut report $(W,J)$, where $J$ is independent uniform on $[d]$, and score $Y=X_J$ after the decision. Every call, including the reissued $W$, is physical and charged. The actual prefix/suffix law is supported on equality of its two $W$ coordinates and admits no nonzero product submeasure. Indeed a nonzero independent product supported on the diagonal must have both $W$ marginals concentrated at one common point, whereas every prefix submeasure has a nonatomic $W$ marginal.

Nevertheless, conditional on $W=w$ the continuation response is the bit vector in $L^2(\mathrm{Unif}[d])$. Theorem~\ref{thm:cutduality} gives exactly the random-access obstruction of Example~\ref{ex:query}; exact prediction still requires $M\ge2^d$, while the final checkpoint needs only two labels. Thus failure of a global certificate neither shows easy compression nor invalidates conditional acquired geometry. No copy of $W$ is retained for free by the algorithm: it is reissued by the specified physical experiment.
\end{example}

\subsection{From continuation charts to a finite recursion}
We give an explicit class where the earlier continuation geometry is both necessary and sufficient in order. Its acquisition is sequential but need not be independent or stationary. Fix $d\ge1$. A prepared experiment first reports $H_1,\ldots,H_d\in\R$ in $d$ paid calls, in that order, with any jointly specified law $\mu$. It next reports a single finite-dimensional packet $U$ by a paid call, and scores a scalar $Y\in[-1,1]$ after the decision. Preparation and the scored audit are also charged. There are exactly $d+3$ raw calls. Report and phase registers are not rereadable across cuts. The only admitted control in this theorem is this acquisition order; the theorem is not an optimization of that order.

Assume that $z(h,u)=\E[Y\mid H=h,U=u]$ has a bounded Borel extension to all $h\in\R^d$ and satisfies
\begin{equation}\label{eq:chartlip}
 |z(h,u)-z(h',u)|\le L\norm{h-h'}_2
 \quad\text{for every }h,h',u.
\end{equation}
Assume $\E(1+|H_j|)^{2+\zeta}\le A$ for some $\zeta>0$, uniformly in $j$. Finally there is a compact Borel set $E\subset\R^d$ with $\mu(E)=p>0$, $\mu|_E$ having Lebesgue density at most $b<\infty$, and a suffix probability $\eta$ such that
\begin{equation}\label{eq:chartlower}
 K(du\mid h)\ge\alpha\eta(du)\ (h\in E),\qquad
 \norm{z(h,\cdot)-z(h',\cdot)}_{L^2(\eta)}
 \ge l\norm{h-h'}_2\ (h,h'\in E),
\end{equation}
where $\alpha,l>0$. These are hypotheses on the law and task response computed from the raw kernels. No covariance, Fisher metric or pre-existing manifold is supplied as geometric input. The second inequality identifies locally distinguishable acquired directions through actual future tests.

\begin{theorem}[Continuation-chart realization]\label{thm:chart}
Under \eqref{eq:chartlip}--\eqref{eq:chartlower}, put
$\xi=(h\mapsto z(h,\cdot))_*(\alpha\mu|_E)$ and let $B$ be the full-history Bayes risk for this one task. There are constants $0<c\le C<\infty$, depending only on $d,\zeta,A,L,p,b,\alpha,l$, such that for every $M\ge1$,
\begin{equation}\label{eq:chartlaw}
 cM^{-2/d}\le q_M(\xi)\le R_\on(M)-B\le CM^{-2/d}.
\end{equation}
The upper uses a genuine $M$-label recursion after each scalar report and after the suffix packet. Its constants do not require independence of the prefix coordinates. If $W_M^{\rm all}$ denotes the maximum of terminal quantization and all the unconditional and conditional cut lower bounds above, then
\begin{equation}\label{eq:chartcomplete}
 R_\on(M)-B\asymp W_M^{\rm all}\asymp q_M(\xi)\asymp M^{-2/d}.
\end{equation}
Thus these cut widths are quantitatively sufficient for this declared sequential-chart class, not for all causal experiments. The phase has $d+2$ pre-audit values; if stored, it multiplies the persistent count. The scalar task always has $R_\cp(M)-B\le M^{-2}$.
\end{theorem}
\begin{proof}
A Hilbert ball of radius $r$ meeting the image of $E$ has preimage contained in a Euclidean ball of radius $2r/l$, by \eqref{eq:chartlower}. Its $\xi$-mass is at most $\alpha b v_d(2r/l)^d$, where $v_d$ is the volume of the Euclidean unit ball. Taking
\[
 r=\frac l2\left(\frac{p}{2b v_d M}\right)^{1/d}
\]
shows that $M$ such balls cover at most half the actual mass $\alpha p$. The remaining mass contributes at least $\alpha p r^2/2$. Theorem~\ref{thm:cutduality} proves the first two inequalities of \eqref{eq:chartlaw}.

For the construction we give the finite scalar quantizer explicitly. Put $a_0=\zeta/2$ and, for integer $k\ge1$, map $x\ge0$ to
\[
 t(x)=1-(1+x)^{-a_0},\quad j=\lfloor kt(x)\rfloor,\quad
 Q_k(x)=(1-j/k)^{-1/a_0}-1;
\]
extend oddly to negative $x$. It has at most $2k-1$ values. Since $0\le Q_k(x)\le x$ and $t'$ decreases,
\begin{equation}\label{eq:compandpoint}
 |x-Q_k(x)|\le\frac{(1+|x|)^{1+a_0}}{a_0k},\qquad
 \E|H_j-Q_k(H_j)|^2\le\frac{A}{a_0^2k^2}.
\end{equation}
This includes the two unbounded outside cells; truncating a Gaussian without paying its tails is unnecessary.

Let $m=\lfloor M^{1/d}\rfloor$ and $k=\lfloor(m+1)/2\rfloor$. On receiving $H_j$ append its quantizer index. After $j$ reports there are at most $m^j\le M$ possible tuples. No actual real prefix is retained. At the packet call compute $z(\widehat h,U)$ from this finite tuple and the current packet, then retain its nearest midpoint in an $M$-cell uniform partition of $[-1,1]$. The old tuple is discarded. The packet and all arithmetic workspace are erased before the next cut. Equation~\eqref{eq:chartlip}, the marginal estimates \eqref{eq:compandpoint}, and Minkowski give root excess at most
\[
 L\sqrt{dA}/(a_0k)+M^{-1}\le C_0M^{-1/d}.
\]
Projection onto the full-history conditional mean proves the upper. This is a finite-label Borel implementation, not yet a finite-bit implementation of arbitrary real coefficients; those resources are discussed below.

Each cut lower is at most the risk of every online predictor, while the particular $\xi$ lower is at least $cM^{-2/d}$. The proved upper therefore yields \eqref{eq:chartcomplete}, with no assumption that an optimal one-cut center list is recursively executable. At the final checkpoint uniform scalar quantization of $z(H,U)\in[-1,1]$ has error at most $M^{-2}$.
\end{proof}

The hypotheses can be checked without a global smooth parameterization. Only one actual positive-mass chart is needed for the lower, while the marginal moment and global response modulus control all acquired tails for the upper. The hard intermediate memory constraint is resolved by coordinatewise acquisition, not by an assumed contractive posterior. Degenerating chart mass or response rank changes the constants and is not hidden in an ambient dimension. The next theorem verifies every ingredient in a genuinely noisy, correlated continuation experiment.
